\section{Complete Symbol and Computational Index}

This index is normative for the notation used in this document. Local dummy variables are marked as
such; repeated glyphs with different subscripts are not interchangeable.

\small
\begin{longtable}{>{\raggedright\arraybackslash}p{0.16\textwidth}
                  >{\raggedright\arraybackslash}p{0.34\textwidth}
                  >{\raggedright\arraybackslash}p{0.40\textwidth}}
\toprule
Symbol & Meaning and unit & Definition or calculation method\\
\midrule
\endfirsthead
\toprule
Symbol & Meaning and unit & Definition or calculation method\\
\midrule
\endhead
\bottomrule
\endfoot
$a,a_0$ & instantaneous and zero-load equilibrium normal spacing; length & $a$ is the primary
normal coordinate. Compute $a_0$ from Eq.~\eqref{eq:chain_equilibrium}, or from the minimum of a
validated replacement energy.\\
$a^\ddagger,a_{\rm inf}$ & normal barrier and lattice inflection spacing; length & Bracket all roots
of $N_AU_\infty'(a)=\sigma A_0$, classify by $U_\infty''$, and use the negative-curvature root as
$a^\ddagger$. Equation~\eqref{eq:normal_instability} gives $a_{\rm inf}$.\\
$s,s_0,\widetilde s,b$ & unwrapped registry, reference registry, within-cell registry, and slip
repeat; length & $s=zb+\widetilde s$; obtain $b$ and $s_0$ from the chosen lattice geometry and use
$0\leq\widetilde s<b$.\\
$z,\mathbb Z,\ddagger$ & integer well label, set of integers, and transition-state marker;
dimensionless & Increment or decrement $z$ when unwrapped $s$ crosses a periodic cell boundary.
The dagger is not multiplication; it marks a saddle or activation barrier.\\
$r,r_e,r_p,r_{ij}$ & pair distances; length & $r_e$ is Eq.~\eqref{eq:re_sigma}; $r_p$ is
Eq.~\eqref{eq:pair_distance}; $r_{ij}=|\bm R_i-\bm R_j|$ in the selected atomic geometry.\\
$\bm d,\bm n,\bm e$ & unit slip direction, plane normal, and uniaxial loading direction;
dimensionless & Normalize the crystallographic vectors and verify
$\bm d\mathbin{\cdot}\bm n=0$.\\
$\bm R_p^-,\bm r_p,\bm R_{ij}$ & lower-row site, cross-row separation vector, and reference lattice
vector; length & Construct directly from the lattice basis and integer cell indices.\\
$m,n,q$ & repulsive exponent, attractive exponent, and generic exponent; dimensionless & Require
$m>n>1$ for the present one-dimensional sums; they are neither masses nor Schmid factors.\\
$C_{m,n}$ & generalized LJ normalization; dimensionless & Evaluate
Eq.~\eqref{eq:generalized_LJ}.\\
$\epsLJ,\sigLJ$ & pair well depth and zero-crossing distance; energy and length & Obtain from the
declared pair potential or independent cohesive calibration; do not fit them to the fatigue loop.\\
$\begin{gathered}v_{m,n},E_{\rm pair}\\E_N,U_\infty\end{gathered}$ & pair, finite-system, finite-chain, and infinite-chain-per-atom
energies; energy & Use Eqs.~\eqref{eq:generalized_LJ}, \eqref{eq:pair_counting},
\eqref{eq:finite_chain_energy}, and \eqref{eq:normal_chain_energy}.\\
$W_{m,n},\Delta W_{m,n}$ & cross-row energy and registry excess; energy & Evaluate the direct sum or
Eq.~\eqref{eq:full_W_mn}; subtract the $s_0$ value for $\Delta W$.\\
$\gamma_{m,n},\gamma_{\rm EAM},\gamma_{\rm GSF}$ & slip or generalized stacking-fault energy
density; energy/area & Divide exact interface-cell excess energy by its represented area; see
Eqs.~\eqref{eq:gamma_definition}, \eqref{eq:EAM_GSF}, and \eqref{eq:full_GSF}.\\
$\psi,\Phi,\Phi_n$ & tilted slip energy, full patch potential, and normal-only potential; energy &
Use Eqs.~\eqref{eq:tilted_energy} and \eqref{eq:full_patch_energy}; set $s=s_0,\tau=0$ for
$\Phi_n$.\\
$\Delta G_{z,\pm}^\ddagger$ & forward/backward activation barrier; energy & Locate the full-energy
minimum and adjacent saddle by bracketed root finding, classify by the exact second derivative,
then use Eq.~\eqref{eq:exact_barrier}.\\
$A_{\rm at},A_0,A_{\rm cell}$ & area per atomic column, correlated patch area, and interface-cell
area; area & Compute atomic areas from the declared lattice cell. Choose $A_0=N_AA_{\rm at}$ from
a physical correlation area, not an FEM element area.\\
$N,N_A$ & finite atom count and equivalent column count; integers & $N_A=A_0/A_{\rm at}$ must agree
with the energy counting convention.\\
$h_{\rm slip}$ & thickness homogenizing a slip discontinuity; length & Specify from the selected
slip band or coarse-graining thickness independently of mesh size.\\
$\sigma,\sigma_{\rm c},\tau,\tau_{\rm int},\tau_{\rm ideal}$ & applied normal stress, normal
instability stress, resolved shear, internal traction, and ideal shear strength; stress & Use
Eqs.~\eqref{eq:normal_instability}, \eqref{eq:Schmid}, \eqref{eq:traction_full}, and
\eqref{eq:ideal_strength}.\\
$F_n,F_s,\mathcal M$ & normal force, slip force, and resolved-shear projection; force, force,
dimensionless & $F_n=\sigma A_0$, $F_s=\tau A_0$,
$\mathcal M=(\bm d\mathbin{\cdot}\bm e)(\bm n\mathbin{\cdot}\bm e)$, and
$\tau=\mathcal M\sigma$.\\
$\delta,\delta_0,\eta$ & normalized registry, reference registry, and separation; dimensionless &
$\delta=s/b$, $\delta_0=s_0/b$, $\eta=a/b$.\\
$p,k,\ell,i,j$ & direct-row, pair-separation, reciprocal-mode, and atom/component indices; integers
& Their ranges are stated at each sum. $i,j$ in $\delta_{ij}$ label components, not atoms.\\
$\nu,\mu,\chi$ & Mellin exponent, Bessel order, and Bessel argument; dimensionless & In the lattice
sum, $\nu=q/2$, $\mu=\nu-1/2$, and $\chi=2\pi|\ell|\eta$. The Bessel argument is $\chi$, not the
well index $z$.\\
$t,u,x,y$ in transform integrals & local dummy integration/Fourier variables; $y$ is also the
local dependent variable in the displayed Bessel ODE & Their dimensions follow the displayed
substitution and they have no material-state meaning.\\
$\begin{gathered}g(x),\widehat g(\ell),I_0\\I_\ell,\mathcal I\end{gathered}$ & local Poisson-transform function, its Fourier
transform, zero/nonzero Fourier integrals, and a generic Bessel integral & These are derivation-only
objects defined immediately before the equations in which they occur; evaluate them by the exact
Gamma/Bessel identities shown there.\\
$\Gamma(\cdot),\zeta(\cdot),I_\mu,K_\mu$ & Gamma, Riemann-zeta, modified Bessel first-kind, and
modified Bessel second-kind functions & Use standard special-function libraries; verify $K_\mu$
against Eq.~\eqref{eq:K_definition} when testing.\\
$\Zsh_\nu,\Bmn_q,\Delta\Bmn_q$ & shifted Epstein--Hurwitz sum, closed exponent sum, and registry
difference; dimensionless & Use Eqs.~\eqref{eq:Z_definition}, \eqref{eq:Bq}, and
\eqref{eq:DeltaBq}. Increase the reciprocal cutoff to tolerance and confirm test points against
the direct sum.\\
$\beta,\gamma$ in Eq.~\eqref{eq:Bessel_generic_integral} & positive local integration parameters &
For that derivation only, $\beta=\eta^2$ and $\gamma=\pi^2\ell^2$; they are not inverse temperature
or stacking-fault energy.\\
$\begin{gathered}\pi,e,\ii,\operatorname{Re}\\|\cdot|,\exp,\ln,\sum\\
\int,\lim,\max\end{gathered}$ & circle constant, natural-logarithm base, imaginary
unit, real part, norm, and standard operators & $\exp x=e^x$ and $\ln$ is its inverse. Sums and
integrals use the displayed limits; $\lim$, $\max$, and absolute values have their standard exact
meanings.\\
$!,\bmod,\sim,\propto,\mathbb R$ & factorial, remainder operation, asymptotic/equivalence marker,
proportionality marker, and real-number set & $j!=\Gamma(j+1)$ for nonnegative integer $j$;
$s\bmod b$ retains only the within-cell registry, whereas unwrapped $s$ retains slip history.\\
\end{longtable}

\begin{longtable}{>{\raggedright\arraybackslash}p{0.16\textwidth}
                  >{\raggedright\arraybackslash}p{0.34\textwidth}
                  >{\raggedright\arraybackslash}p{0.40\textwidth}}
\toprule
Dynamic symbol & Meaning and unit & Definition or calculation method\\
\midrule
\endfirsthead
\toprule
Dynamic symbol & Meaning and unit & Definition or calculation method\\
\midrule
\endhead
\bottomrule
\endfoot
$t,f,\omega_{\rm L}$ & time, cyclic frequency, and loading angular frequency; time, 1/time, 1/time &
$\omega_{\rm L}=2\pi f$; for nonsinusoidal loading test every relevant harmonic.\\
$T,\kB,\beta_T$ & absolute temperature, Boltzmann constant, and inverse thermal energy; K,
energy/K, 1/energy & $\beta_T=(\kB T)^{-1}$.\\
$q_i,v_i,m_i^*$ & collective coordinate, velocity, and effective mass; length, length/time, mass &
$q_a=a,q_s=s$. For two rigid patches of masses $M_U,M_L$, use
$m_i^*=M_UM_L/(M_U+M_L)$; use the moving mass when the other patch is fixed.\\
$\mathcal K_i,R_i,\tau_{K_i}$ & memory-friction kernel, projected random force, and kernel decay
time & Obtain $\mathcal K_i(t)=\beta_T\langle R_i(0)R_i(t)\rangle$ from a constrained equilibrium
MD projection and estimate $\tau_{K_i}$ from its decay.\\
$\Gamma_i,M_i,D_i$ & Markov friction, long-time mobility, and configuration diffusivity;
mass/time, length$^2$/(energy time), length$^2$/time & If a Markov plateau exists,
$\Gamma_i=\int_0^\infty\mathcal K_i(t)\dd t$, $M_i=1/\Gamma_i$, and $D_i=\kB TM_i$.\\
$\xi_i,\delta_{ij},\delta(t-t')$ & unit white noise, Kronecker delta, and Dirac delta & Normalize by
$\langle\xi_i(t)\xi_j(t')\rangle=\delta_{ij}\delta(t-t')$. Do not confuse $\delta_{ij}$ with
$\delta=s/b$.\\
$\rho,\rho_*,\mathcal P_b,\dd\Lambda$ & phase density, reference phase density, bonded phase domain,
and phase-volume element & Normalize $\rho$ on $\mathcal P_b=\Omega_b\times\mathbb R^2$;
$\rho_*$ only nondimensionalizes the logarithm.\\
$\mathcal J_{v_i}^{\rm irr}$ & irreversible velocity-space probability current & Compute
Eq.~\eqref{eq:velocity_irreversible_current}; it produces the Markov Kramers dissipation.\\
$\mathcal F_K,\dot D_K$ & inertial nonequilibrium free energy and irreversible loss rate; energy,
energy/time & Use Eqs.~\eqref{eq:Kramers_free_energy} and \eqref{eq:Kramers_dissipation}.\\
$\epsilon_i^{\rm in},\epsilon_i^{\rm mem}$ & inertial and memory reduction tests; dimensionless &
Compute Eq.~\eqref{eq:overdamped_conditions}; use Smoluchowski only when both are small for every
relevant loading harmonic.\\
$P,P_{\rm eq},P_*,\mathcal Z_{\rm conf}$ & configuration density, equilibrium density, reference
density, and configurational partition function & $P$ is the velocity marginal of $\rho$; normalize
using Eq.~\eqref{eq:normalization}; use $P_{\rm eq}$ only in a bounded bonded basin.\\
$\begin{gathered}\Omega_b,\partial\Omega_b\\\partial\Omega_{\rm fr},\bm n_\Omega,\dd S\end{gathered}$ & bonded domain, its
boundary, absorbing fracture portion, outward normal, and boundary measure & Place the normal
fracture boundary at the full-energy saddle $a^\ddagger$, not at a fitted distance.\\
$J_a,J_s,\bm J$ & configuration probability currents; 1/(length time) & Evaluate
Eqs.~\eqref{eq:Ja}--\eqref{eq:Js}; use conservative finite-volume face fluxes.\\
$\mathcal F,\mu_P,D_{\rm irr},\dot D_{\rm irr}$ & overdamped free energy, probability chemical
potential, accumulated dissipation, and rate; energy, energy, energy, energy/time & Use
Eqs.~\eqref{eq:free_energy_functional}, \eqref{eq:probability_chemical_potential},
\eqref{eq:dissipation}, and \eqref{eq:accumulated_dissipation}.\\
$S,P_{\rm crack},p_{\rm tail}$ & intact survival, cumulative crack probability, and instantaneous
tail probability; dimensionless & Use Eqs.~\eqref{eq:survival_flux} and
\eqref{eq:tail_probability}; never identify dissipation directly with crack probability.\\
$\bar a,\bar s,\operatorname{Var}[a],\bar E$ & unconditioned means, variance, and mean potential
energy & Use Eqs.~\eqref{eq:mean_a}--\eqref{eq:mean_energy}; divide by $S$ for intact-conditional
moments.\\
$\begin{gathered}g,\bm\nabla,\langle\cdot\rangle\\\partial_t,\dot{(\cdot)},(\cdot)'\end{gathered}$ & generic observable, gradient,
ensemble average, partial time derivative, total time derivative, and one-variable derivative &
Use Eq.~\eqref{eq:general_moment_boundary}; keep its boundary term for absorbing fracture.\\
\end{longtable}

\begin{longtable}{>{\raggedright\arraybackslash}p{0.16\textwidth}
                  >{\raggedright\arraybackslash}p{0.34\textwidth}
                  >{\raggedright\arraybackslash}p{0.40\textwidth}}
\toprule
Slip/EAM symbol & Meaning and unit & Definition or calculation method\\
\midrule
\endfirsthead
\toprule
Slip/EAM symbol & Meaning and unit & Definition or calculation method\\
\midrule
\endhead
\bottomrule
\endfoot
$s_z,s_{z,\pm}^\ddagger$ & stable registry in well $z$ and adjacent forward/backward saddle; length
& Solve $\partial_s\gamma=\tau$ on each periodic interval and classify with
$\partial_s^2\gamma$.\\
$\gamma_p,\varepsilon_p$ & reduced plastic shear and projected axial plastic strain; dimensionless
& Use Eq.~\eqref{eq:plastic_strain}; this is ideal single-slip bookkeeping, not hardening.\\
$s_L,s_R,x,y,T_{\rm FP}$ & reflecting boundary, absorbing boundary, quadrature variables, and mean
first-passage time; length, length, length, length, time & Select boundaries around the exact well
and saddle and evaluate Eq.~\eqref{eq:exact_MFPT} by adaptive quadrature.\\
$k_{z,\pm}^{\rm FP},k_{z,\pm}$ & first-passage and optional asymptotic transition rates; 1/time &
Invert the exact overdamped MFPT under its assumptions; use Eq.~\eqref{eq:Kramers_rate} only for a
verified high barrier.\\
$p_z,J_{a,z}$ & probability density and normal current in well $z$ & Advance
Eq.~\eqref{eq:master_equation} only when intrawell, escape, and load time scales separate.\\
$\begin{gathered}E_{\rm EAM},\phi,F,\rho_i\\f,f_0,\beta_\rho\end{gathered}$ & EAM energy, pair term, embedding function, host
density, density kernel, amplitude, and decay coefficient & Use a complete validated EAM
parameterization. Sum $f(r_{ij})$ before applying nonlinear $F$. $\beta_\rho$ is not $\beta_T$.\\
$\mathcal U,\mathcal L,\mathcal U_{\rm cell}$ & upper half-crystal, lower half-crystal, and upper
interface cell & Generate from the selected crystal orientation and converge neighbor extent or the
validated EAM cutoff.\\
$M_U,M_L$ & masses of upper and lower representative moving patches; mass & Sum atomic masses in
the correlated patch used for the relative collective coordinate.\\
\end{longtable}
\normalsize

\subsection{Reproducible calculation sequence}

\begin{enumerate}
  \item Declare geometry and counting: $m,n,\epsLJ,\sigLJ,b,A_{\rm at},A_0,s_0$ and whether the
  normal-only or optional slip branch is active.
  \item Evaluate $U_\infty$ and $W_{m,n}$. Increase the Bessel-mode cutoff until the energy and all
  required derivatives are stable at the requested tolerance; verify representative points against
  the direct real-space sum.
  \item Locate minima, saddles, and instability points by bracketed root finding on the full
  analytic energy and classify them with full derivatives. No polynomial energy fit is required.
  \item Obtain $m_i^*$ and $\mathcal K_i$ from atomistic geometry and constrained equilibrium MD.
  Test Eq.~\eqref{eq:overdamped_conditions}; do not introduce $M_i$ unless the reduction is
  supported.
  \item If inertia matters, advance Eq.~\eqref{eq:Kramers_PDE} with a conservative,
  positivity-preserving finite-volume method in $(q_i,v_i)$. If the reduction is justified, advance
  Eq.~\eqref{eq:Smoluchowski} with conservative face fluxes. Check positivity, normalization, and
  time-step and domain convergence in either case.
  \item For crack initiation, place the absorbing normal boundary at the full-energy saddle,
  integrate its outward probability flux, and report $p_{\rm tail}$, $S$, and $P_{\rm crack}$
  separately. Use Eq.~\eqref{eq:dissipation_with_boundary}, not the closed-boundary balance.
  \item Repeat with finer state grids, smaller time steps, larger lattice-sum cutoffs, and larger
  velocity domains. Report convergence before interpreting a fatigue hysteresis loop.
\end{enumerate}
